\subsection{Hierarchical Anomaly Detection and Fault-Domain Diagnosis}
\label{sec:hierarchical_diagnosis}

HS-AeroTS separates runtime anomaly detection from conditional fault-domain diagnosis according to Equation~\eqref{eq:problem_cascade}. Its input is the 1566-dimensional feature vector $\mathbf{f}_w$ defined in Section~\ref{sec:data_preprocessing}, and its outputs are a Stage 1 anomaly score, a Stage 2 fault-domain probability vector when diagnosis is performed, and a final runtime class under the Cascade. Stage 1 is a class-balanced binary LightGBM model configured with 1600 estimators, learning rate 0.025, 64 leaves, minimum child size 80, subsample 0.90, feature subsampling 0.75, $\ell_1$ regularization 0.20, $\ell_2$ regularization 8.0, and 80-round early stopping. It produces an anomaly score $p_1(\mathbf{f}_w)\in[0,1]$ for each feature window. Its operating threshold $\tau$ is selected by the implemented precision--recall search on validation labels, using the threshold with maximum F1, and is then frozen for test evaluation. AUPRC and AUROC assess threshold-independent detection, while threshold-dependent F1 and event-level measures assess the operating point.

Stage 2 is a four-class LightGBM model configured with 1200 estimators, learning rate 0.025, 64 leaves, minimum child size 30, subsample 0.90, feature subsampling 0.75, $\ell_1$ regularization 0.20, $\ell_2$ regularization 8.0, class-balanced weighting, and 80-round early stopping. It takes the same feature vector and is trained only on windows with an explicit anomalous fault-domain label $y_w^{\mathrm{dom}}\in\{1,2,3,4\}$. These codes denote External Position, Global Position, Altitude, and Mechanical/Electrical, respectively. Normal windows have code 0 and Uncategorized windows have code 5; both are excluded from this conditional training task. Its output is a four-class probability vector $\mathbf{p}_2(\mathbf{f}_w)$ and a predicted domain $\hat{y}_w^{\mathrm{dom}}=\arg\max_d p_{2,d}(\mathbf{f}_w)$. Stage 2 uses the same log membership as Stage 1 and is evaluated with Macro-F1, Balanced Accuracy, Macro Recall, and per-class metrics. Core model parameters are reported in Table~\ref{tab:protocol}, while frozen machine-readable configurations and run records are part of the analysis package described in the Data Availability Statement.

Thus, Stage 1 is a binary screening task, Stage 2 is a conditional four-class diagnosis task, and Cascade is the complete five-class runtime decision formed by gating Stage 2 with Stage 1. Direct Five-Class is the comparator that predicts Normal plus the four fault domains in one five-class LightGBM model without this gate. Random Forest uses the configured 500 trees, square-root feature subsampling, minimum leaf size 2, and balanced subsampling; majority and stratified-random methods are used as reference methods under the corresponding task definitions. The distinction is important for interpreting results: conditional Stage 2 metrics are not complete five-class Cascade metrics, and the Cascade is not assumed to outperform Direct Five-Class under every protocol.
